\chapter{Conclusion}

This internship provided a valuable opportunity to design, build, and finalize a cross-platform computational application for compressible flow analysis. Through software engineering principles, numerical analysis, and algorithmic optimization, a production-ready aerodynamics calculator was successfully developed for mobile platforms.

A comprehensive suite of solvers was successfully engineered using the Flutter framework and Dart language, and their performance was evaluated. Key computational strategies, such as Newton--Raphson iterations for implicit isentropic relations, bisection methods for friction and heat addition boundaries, analytical cubic root-finding for oblique shock angles, and adaptive RK45 integration for conical shocks, were successfully implemented. A decoupled three-layer architecture and a custom expression parser were also developed, enabling an interactive user interface and a clean computation layer.

The overall project demonstrated that cross-platform application development is highly effective for delivering complex engineering tools when the algorithmic logic is optimized and decoupled. The successful testing and deployment pipeline verify the readiness of the application for public release on the Google Play Store.

Overall, this internship enhanced my knowledge of numerical methods, cross-platform app design, and software testing. It also improved my problem-solving skills, analytical thinking, and understanding of the software development lifecycle, making it a valuable learning experience.
